% =====================================================================
%  28.5  Recognition and the method boundary -- the shape signature
%  (label-free, order-free, whole-shape object, static tree), and the
%  hand-offs: values/colours must match -> Ch.13 DP / Ch.17 DSU-on-tree;
%  order matters -> ordered encodings (the CF 1671/E "Preorder" family,
%  NAME ONLY -- its URL lives in this chapter's mixed set); path strings
%  with labels -> suffix machinery; online edits -> Ch.22 HLD / Ch.25 LCT;
%  distances -> Ch.12 metrics / Ch.24 centroid; general graphs -> colour
%  refinement, not AHU.
%  WORKED pick examples only -- no solvedproblem (contract: two, spent on
%  CF 601D and CF 1800G).
%  Exports: sec:c22-boundary  def:c22-signature  def:c22-pick
%  lst:c22-kit  fig:c22-handoff
%  CF 1182/D "Complete Mirror" (ch20_advanced_tree_dp), CF 914/E
%  "Palindromes in a Tree" (ch18_centroid_decomposition) and CF 375D /
%  600E / 1009F (Ch.17) are PROSE references only -- their URLs live in
%  other chapters and nothing here links a foreign URL.
% =====================================================================
\section{The Method Boundary: When Shape Is Not the Question}\label{sec:c22-boundary}

\begin{previouslybox}[Chapter 27 handed the question over; this section returns the rest]
\(\secref{c21-boundary}\) closed with a debt: ``which subtree equals which'' is not a
counting question and Mo cannot answer it -- Chapter~28 owns it.  Collected, four
tools now sit in the kit: the interned AHU labeller (\(\secref{c22-ahu}\)), the
commutative subtree hash and the positional string hash (\(\secref{c22-hash}\)), the
centre-rooted canonical form \(\Phi\) (\(\secref{c22-centre}\)), and the id census
(\(\secref{c22-distinct}\)).  A tool this blunt reads \emph{structure only}, so it is
wrong in a specific, nameable way: whenever a \emph{name}, a \emph{value}, an
\emph{order}, a \emph{distance} or an \emph{edit} carries information, this chapter's
machinery quietly throws it away.  The boundary here is not about difficulty -- a
rejected problem is usually easier than an accepted one -- it is about which
information the statement actually uses.
\end{previouslybox}

\begin{topicintro}
Ask four questions in order.  \emph{Does renaming every vertex leave the answer
alone?}  No \(\Rightarrow\) the objects are \emph{labelled} trees: this chapter's
ids and hashes are blind to names, so the job is a value fold
(\chref{c10}) or a keyed merge (\chref{c28}).  \emph{Does reordering a vertex's
children leave the answer alone?}  No \(\Rightarrow\) the object is an
\emph{ordered} tree: skip the sort and use an in-order / preorder encoding, where
suffix tools, not shape hashes, are the right algebra.  \emph{Is the queried object
a whole downward shape (subtree, component) rather than a path or a distance?}  No
\(\Rightarrow\) Ch.~12's metrics (\chref{c09}) or centroid
(\chref{c18}).  \emph{Is the tree fixed while the queries run?}  No \(\Rightarrow\)
canonical forms must be rebuilt after every edit: HLD (\chref{c16}) and the
link-cut tree (\chref{c23}) hold dynamic structure; a static
\(\bigO{n\log n}\) fold does not compete.  Pass all four and the problem is
this chapter's, and \defref{c22-pick} says which of the four tools to fire first.

\textbf{Prerequisites.} The whole chapter, plus the boundary stations of
Ch.~13/17/22/24/25/27 that take the work back.

\textbf{After this section you can}
\begin{itemize}[nosep]
  \item recite the shape signature: label-free, order-free, whole-shape object, static tree; \hfill\lv{4}
  \item reject isomorphism tooling for value-bearing, order-bearing, distance-bearing and dynamic problems; \hfill\lv{5}
  \item choose id-census vs.\ commutative hash vs.\ positional string hash by what is being counted; \hfill\lv{5}
  \item say what AHU cannot do for general graphs, and what replaces it there. \hfill\lv{4}
\end{itemize}
\end{topicintro}

% ---------------------------------------------------------------
\subsection{The signature}
% ---------------------------------------------------------------

\begin{definition}{The shape signature (four tests)}{c22-signature}
A problem is isomorphism-shaped iff all four hold:
\begin{enumerate}[nosep,leftmargin=1.6em]
  \item \textbf{Label-free:} a bijection of vertex names must not change the
        answer.  Vertex \emph{values} that must agree are legal only as extra
        payload inside the interned tuple (\(\secref{c22-distinct}\)'s third
        hardening item); values that get \emph{combined} are not this chapter.
  \item \textbf{Order-free:} permuting any vertex's child list must not change
        the answer.  This is the test that separates \(\defref{c22-shapehash}\)
        (commutative) from \(\defref{c22-strhash}\) (positional).
  \item \textbf{Whole-shape object:} the compared thing is an entire downward
        subtree or a component left by a cut -- something a bottom-up fold
        captures in one key.  Paths, distance profiles and ancestor sets are
        not such objects.
  \item \textbf{Static tree:} the edges do not change while questions arrive, so
        one \(\bigO{n\log n}\) canonicalisation amortises over all compares.
\end{enumerate}
Fail any test and the table below names where the problem really lives.
\end{definition}

\begin{recognitionbox}[Question shapes: this chapter versus each look-alike]
\reckey{statements that compare or count \emph{structures}: ``are these two trees the
same tree'', ``identical'', ``the same structure'', ``can the children be ordered so
that \(\dots\)'', ``all pieces the same after one cut'', ``how many different shapes
occur''.  The tell is that renaming every vertex and permuting every child list
cannot change the answer -- tests 1--2 of \(\defref{c22-signature}\) -- and that the
thing being compared is an entire downward shape, on a tree nobody edits.}
\recuse{
  \item two trees, compared once -- AHU plus the centre form \(\Phi\), \(\bigO{n\log n}\)
        and then \(O(1)\) per compare (\(\secref{c22-ahu}\), \(\secref{c22-centre}\),
        \(\coderef{c22-unrooted}\));
  \item ``can the children mirror-pair?'' at every vertex -- shape ids with the
        one-odd-class test (\(\thmref{c22-sym}\), solved CF 1800G);
  \item how many shapes / who shares my shape -- the id census
        (\(\thmref{c22-classes}\), \(\coderef{c22-classes}\));
  \item thousands of subtree-vs-subtree compares inside one tree -- the commutative
        hash (\(\defref{c22-shapehash}\));
  \item distinct downward \emph{strings} per node -- positional hashes merged
        small-to-large (\(\defref{c22-strhash}\), solved CF 601D);
  \item every component of one cut identical -- hashes plus the reroot view and a
        per-candidate \(\Phi\) check (\(\probref{c22-summit}\)).
}
\recnotuse{
  \item vertex \emph{values} are combined (sums, xors, minima) -- test~1 fails: tree
        DP (\(\chref{c10}\)) or a keyed per-node merge (\(\chref{c28}\); CF 375D,
        600E, 1009F are treated there).  If values must \emph{match} rather than
        combine, intern \((\text{child ids},\text{value})\) jointly and the shape
        id becomes one component of the key, never the whole of it;
  \item child order is semantic -- left/right roles, preorder strings, ``first child
        then second'': test~2 fails, so the commutative fold is a \emph{wrong
        equivalence} (two different strings, one hash, false YES with probability
        \(1\)).  Use the traversal string or the positional hash; the perfect
        binary tree preorder family is the CF 1671/E ``Preorder'' shape, named here
        only and waiting in the mixed set;
  \item the object is a path or a distance -- test~3 fails: Ch.~12 metrics
        (\(\chref{c09}\)), and for pair distributions or path-shaped palindromes the
        centroid machinery (\(\chref{c18}\), where CF 914/E is solved);
  \item edges are cut and linked while questions arrive -- test~4 fails: HLD
        (\(\chref{c16}\)) or the link-cut tree (\(\chref{c23}\));
  \item the input is a general graph -- there is no postorder to fold, so AHU does
        not run; colour refinement (the \(k\)-WL test) is the honest fallback.
}
One-line test: \emph{strip the labels, shuffle the children, and check that nothing
changes -- then compare the shape, taking the id for certainty and the hash for
speed.}  If any information survives the stripping, that information is the problem,
and it belongs to a neighbour chapter.
\end{recognitionbox}

% ---------------------------------------------------------------
\subsection{The pick, worked: four specs, four verdicts}
% ---------------------------------------------------------------

\begin{figure}[H]
\centering
\begin{tikzpicture}[diagram, scale=0.95, transform shape,
    mnode/.style={circle, draw=#1, line width=0.9pt, fill=#1!8, minimum size=10mm,
                  font=\scriptsize\bfseries, text=#1!60!black}]
  \node[mnode=primaryblue] (q) at (0,0) {tests};
  \node[mnode=successgreen] (a) at (4.4,1.5)  {1 pass};
  \node[mnode=difforange]   (b) at (4.4,0.5)  {2 pass};
  \node[mnode=softpurple]   (c) at (4.4,-0.5) {3 pass};
  \node[mnode=warningred]   (d) at (4.4,-1.5) {4 pass};
  \node[annot, font=\scriptsize, anchor=west, text width=6.3cm, align=left] at (6.0,1.5) {values matter: fold them (Ch.13) or key them (Ch.17)};
  \node[annot, font=\scriptsize, anchor=west, text width=6.3cm, align=left] at (6.0,0.5) {order matters: positional hash, traversal encoding, suffix tools};
  \node[annot, font=\scriptsize, anchor=west, text width=6.3cm, align=left] at (6.0,-0.5) {distances matter: Ch.12 metrics, centroid for pair distributions};
  \node[annot, font=\scriptsize, anchor=west, text width=6.3cm, align=left] at (6.0,-1.5) {topology changes: HLD, link-cut; re-canonicalising per edit is fatal};
  \node[annot, font=\scriptsize, anchor=west, text width=4.6cm, align=left] at (-3.4,0) {all four pass: canonical id for certainty, commutative hash for speed, centre form for unrooted, census for counting};
  \draw[flow] (q) -- (a); \draw[flow] (q) -- (b);
  \draw[flow] (q) -- (c); \draw[flow] (q) -- (d);
\end{tikzpicture}
\caption{The four failed tests and their destinations.  Each arrow is a
\emph{cheaper} answer, not a harder one: the chapter's tools are the residue after
values, order, distance and dynamism have all been ruled out.}
\label{fig:c22-handoff}
\end{figure}

\begin{center}\renewcommand{\arraystretch}{1.25}\setlength{\tabcolsep}{3pt}
\begin{tabular}{@{}p{0.20\linewidth} p{0.24\linewidth} p{0.26\linewidth} >{\raggedright\arraybackslash}p{0.16\linewidth}@{}}
\toprule
\textbf{Axis} & \textbf{This chapter} & \textbf{The neighbour} & \textbf{Wins}\\
\midrule
Object & a whole shape (subtree, component) & a path, a distance, a value sum & neighbour \\
Key & interned id / commutative hash & prefix value / bitmask & tie (different keys) \\
Certainty & id: exact; hash: \(\le\) depth\(/\,2^{64}\) wrong & exact arithmetic & id \\
Cost per compare & \(O(1)\) after \(\bigO{n\log n}\) & re-fold \(\bigO{\size}\) & this chapter \\
Edits & none (static) & \(O(\log n)\) per op (HLD, LCT) & neighbour \\
Order sensitivity & none by construction & native (encoding, string) & neighbour \\
\bottomrule
\end{tabular}
\end{center}

\begin{tracebox}[Four specs, run through the signature]
\textbf{(a)} \(n\le2\cdot10^5\); for every vertex, does some other vertex have an
isomorphic subtree?  Tests 1--4 pass; the object is a downward shape.
\(\Rightarrow\) \emph{this chapter}: one labelling, then test
\(m(\mathrm{cl}(v))\ge2\) (\(\thmref{c22-classes}\)).  Cost \(\bigO{n\log n}\) total.
\textbf{(b)} Same tree, each vertex carries a value; ask for two vertices with
isomorphic subtrees \emph{and} equal subtree sums.  Test~1 fails on the sum (a
value gets combined, not matched).  \(\Rightarrow\) \emph{not} a bare id census:
intern \((\text{child ids}, \text{value})\) if the value must \emph{match}, else
key by \((\mathrm{cl}(v), \mathit{sum}(v))\) after Ch.~13's fold -- the shape id
becomes one component of the key, never the whole of it.
\textbf{(c)} Perfect binary tree on \(2^k-1\) vertices, letters on vertices; is the
preorder string of \(x\) equal to that of \(y\)?  Test~2 fails (left and right
children are named roles), so the commutative fold is illegal and the
\emph{positional} combine is correct: hash the traversal itself, the
\(\mathit{val}(x_i)B^{\depth(x_i)}\) discipline of \(\defref{c22-strhash}\).  This
is the CF 1671/E ``Preorder'' shape -- by name here, its own solve is a mixed-set
entry.
\textbf{(d)} \(n,q\le10^5\); alternate ``is the tree rooted at \(r\) isomorphic to
the tree rooted at \(r'\)?'' queries with interleaved edge cuts and links.  Test~4
fails.  \(\Rightarrow\) dynamic structure: the link-cut view (\chref{c23}) keeps the
forest; re-running \(\Phi\) after every edit is
\(\bigO{q\cdot n\log n}\) and \(\tle\).  \(\Box\)
\end{tracebox}

\begin{definition}{The order of the four questions (pick protocol)}{c22-pick}
\emph{Do names or values carry information?} \(\Rightarrow\) tree DP
(\chref{c10}) or DSU-on-tree (\chref{c28}); add the shape id to the key only if the
statement still matches shapes.  \emph{Does child order carry information?}
\(\Rightarrow\) an ordered encoding (preorder/parenthesis string, positional hash):
compare as \emph{strings}, and hand long common-prefix questions to the suffix
machinery.  \emph{Is the object a path or a distance?} \(\Rightarrow\) Ch.~12
(\chref{c09}); all-pairs distributions \(\to\) centroid (\chref{c18}).  \emph{Does
the topology change?} \(\Rightarrow\) HLD (\chref{c16}) / link-cut (\chref{c23}).
\emph{Still here?} The problem is shape-shaped, and the second decision is
certainty versus speed: two trees to compare once \(\to\) one AHU pass each and the
centre form \(\Phi\) (\(\secref{c22-ahu}\), \(\secref{c22-centre}\)); thousands of
compares inside one tree \(\to\) commutative hashes (\(\secref{c22-hash}\)); many
distinct objects per node \(\to\) hashes plus small-to-large (CF 601D); ``how many
shapes / who shares mine'' \(\to\) the id census (\(\secref{c22-distinct}\)); ``all
components of one cut identical'' \(\to\) hashes \(+\) the reroot view
(\(\probref{c22-summit}\)).  \emph{And on general graphs?} Nothing here: tree
isomorphism is linear-time precisely because of the bottom-up fold; graph
isomorphism is not known to be, and the honest tool there is colour refinement
(the \(k\)-WL test) plus special structure -- AHU on a graph with a cycle has no
postorder to run.
\end{definition}

\begin{lstlisting}[caption={The pick protocol as a decision skeleton -- paste-ready comments for contest code; each branch names the chapter that owns the technique}, label={lst:c22-kit}]
// 1. do vertex VALUES or NAMES have to agree or combine?   -> NO shape-only key.
//      combine -> tree DP (Ch.13, c10); per-node keyed merge -> DSU on tree
//      (Ch.17, c28). match-and-shape -> intern (child ids, value) together.
// 2. does CHILD ORDER matter (left/right, preorder string)? -> NO commutative fold.
//      use the positional/depth-weighted hash or the traversal string itself;
//      long shared prefixes -> suffix tools.  (CF 1671/E "Preorder" family.)
// 3. is the object a PATH or a DISTANCE, not a downward shape?
//      metrics -> Ch.12 (c09); pair distributions / path palindromes
//      -> centroid (Ch.24, c18).  (CF 914/E is treated there.)
// 4. do edges get cut/linked between questions?            -> NO static canonical form.
//      paths with updates -> HLD (Ch.22, c16); dynamic forest -> LCT (Ch.25, c23).
// 5. STILL HERE -- the object is a SHAPE:
//      two trees, once          -> AHU + centre form Phi (28.1, 28.3; lst:c22-unrooted)
//      many compares, one tree  -> commutative hash, O(1) each (28.2)
//      YES/NO needs certainty   -> interned ids, never hashes (28.1)
//      hash says YES and it matters -> re-verify by explicit bijection DFS (28.2)
//      how many shapes / class sizes -> id census + freq[] (28.4; lst:c22-classes)
//      distinct downward STRINGS      -> positional hash + small-to-large (CF 601D)
//      every component of one cut     -> all hashes + reroot + per-candidate Phi (Summit)
// 6. general graph, not a tree? -> stop: no postorder, use colour refinement (k-WL).
\end{lstlisting}

\begin{insightbox}[The first decision is the object, the second is certainty]
Two questions are worth separating, because students answer them in the wrong
order.  \emph{Is the compared object a shape?} -- decided by the four tests, and
answered by the statement, not by your taste.  \emph{Id or hash?} -- a
\emph{cost} decision inside the chapter: ids are exact but need one shared table
and \(\bigO{n\log n}\) per canonicalisation, hashes are \(O(1)\) per compare but
charge a collision premium (\(\thmref{c22-collision}\)).  When the judge prints a
YES/NO on a single pair, take the exact route -- the collision ledger is a
\emph{budget}, not a licence.  When the batch is thousands of compares inside one
tree, the exact route costs \(\bigO{n}\) per compare and the hash route costs
\(O(1)\): pay the premium.  Everything else in this chapter is a special case of
those two sentences.
\end{insightbox}

\begin{warningbox}[Three boundary errors]
\begin{enumerate}[nosep,leftmargin=1.6em]
  \item \textbf{Shape-keying a value question.} ``Two vertices whose subtrees are
        isomorphic \emph{and} whose values sum to the same'' is not the census of
        \(\secref{c22-distinct}\): the id alone has thrown the values away.
        Symptom: the answer counts pairs that plainly differ in payload.
  \item \textbf{Commutative keying an order question.} A fold that sorts (or adds)
        child keys identifies \(\langle A,B\rangle\) with \(\langle B,A\rangle\).
        For a preorder string that is a \emph{wrong equivalence}: two different
        strings, one hash, false YES with probability \(1\), not \(2^{-64}\).
  \item \textbf{Re-canonicalising under edits.} A cut changes the centre, the
        postorder and every id above the cut.  Reusing stale ids after an edit is
        the classic silent WA -- and rebuilding \(\Phi\) per edit is the classic
        \(\tle\).  Both are the same mistake: test~4 ignored.
\end{enumerate}
\end{warningbox}

\begin{notebox}[Boundary homework this chapter leaves elsewhere]
Three named cousins keep walking out of these doors.  CF 1182/D ``Complete
Mirror'' asks the unrooted version of 1800G's question and is solved in
Ch.~20's advanced tree DP (its URL lives there, so it is named, not linked).
CF 914/E ``Palindromes in a Tree'' asks for path-shaped objects -- centre
arguments in a different metric -- and belongs to centroid (\chref{c18}), where
it is treated.  CF 375D / 600E / 1009F are the per-node keyed folds this chapter
hands to DSU-on-tree (\chref{c28}).  And CF 1671/E ``Preorder'' -- order-bearing,
test~2 failing -- waits in the mixed set below as drill, not as a solved problem.
\end{notebox}

\begin{complexitybox}[The boundary, priced]
Shape-shaped: AHU / \(\Phi\) \(\bigO{n\log n}\) per canonicalisation, \(O(1)\)
per compare afterwards; commutative hash \(\bigO{n}\) precompute, \(O(1)\) per
compare, collision premium \(\le \depth/2^{64}\); id census \(\bigO{n}\) extra
space.  Value-shaped: Ch.~13 fold \(\bigO{n}\) but \(\bigO{\size}\) per subtree
re-read; keyed merge \(\bigO{n\log n}\).  Order-shaped: traversal-string hashing
\(\bigO{n}\) with \(O(1)\) compares, suffix tools \(\bigO{n}\) build.
Distance-shaped: Ch.~12 \(\bigO{n}\) per source, centroid \(\bigO{n\log n}\) for
pair distributions.  Dynamic: HLD \(\bigO{\log^2 n}\), link-cut \(\bigO{\log n}\)
amortised per operation.  The tell is that the \emph{wrong} choice is not slower
by a constant -- it is wrong by a whole information channel, which is why the
protocol order in \defref{c22-pick} comes before any optimisation.
\end{complexitybox}

\begin{recallbox}[28.5 in one breath]
\emph{Four tests make a shape problem: label-free, order-free, whole-shape object,
static tree (\(\defref{c22-signature}\)).  Each failed test is a hand-off, not a
dead end: values \(\to\) Ch.13/Ch.17; order \(\to\) traversal encodings and suffix
tools; paths and distances \(\to\) Ch.12/Ch.24; edits \(\to\) Ch.22/Ch.25; general
graphs \(\to\) colour refinement.  Ask the questions in this order
(\(\defref{c22-pick}\), \(\coderef{c22-kit}\)), and only once the object is
confirmed to be a shape do you choose between the exact id and the fast hash.}
\selftest{
  \item Which test does a preorder-string equality problem fail, and what replaces the commutative fold? \emph{(Test 2, order-free: use the positional/depth-weighted hash or the traversal string, since sorting children identifies different strings.)}
  \item Values must match AND shapes must match: name the one-line change to the interned tuple. \emph{(Intern the sorted tuple of child ids together with the vertex value -- the key becomes (child ids, value), so classes respect both.)}
  \item Why is re-running the centre form after every edge edit fatal, and what is the honest destination? \emph{(Canonicalisation is O(n log n) per call and the centre/postorder change with every edit; dynamic structure queries go to HLD (Ch.22) or link-cut (Ch.25).)}
  \item What makes tree isomorphism linear-ish while graph isomorphism is open territory? \emph{(The bottom-up fold: a tree has a postorder and rooted subtrees decompose; a graph has neither, so the fallback is colour refinement (k-WL).)}
}
\end{recallbox}
